
\documentclass[11pt, a4paper]{article}
\usepackage[utf-8]{inputenc}
\usepackage[margin=1in]{geometry}
\usepackage{graphicx}
\usepackage{hyperref}
\usepackage{amsmath}
\usepackage{amssymb}
\usepackage[numbers]{natbib}
\usepackage{setspace}

\onehalfspacing

\title{AI Tools and Founder Decision-Making Speed and Quality in Early-Stage Startups: A Systematic Literature Review}
\author{Automatic Research Machine}
\date{\today}

\begin{document}

\maketitle

\begin{abstract}
**Background** Artificial intelligence tools are increasingly embedded in early-stage startups, shaping how founders gather information and commit resources under uncertainty. Yet evidence on AI-enabled founder decision-making remains fragmented across disciplines and methodologies, impeding a comprehensive account of its effects.

**Objective** This review synthesizes evidence on how AI tools affect founder decision-making speed and quality in early-stage startups, using the Technology Acceptance Model as its primary theoretical lens.

**Methods** A systematic search of Semantic Scholar, arXiv, and Crossref identified 43 papers, of which 39 were retained for synthesis following PRISMA guidelines after deduplication and screening. Retained studies span regression, survey, case study, and meta-analytic designs published predominantly within the last three years.

**Results** Across the 39 included studies, AI tools are consistently associated with improved decision-making precision, more effective opportunity recognition, and enhanced idea generation among founders. AI capabilities are also linked to gains in innovation, risk mitigation, and competitive advantage, although several studies note that these benefits are constrained by AI's limited capacity for human judgment, emotion, and creativity in certain decisions. One study reports a moderate positive correlation between a generative AI tool's perceived accuracy and relevance (r = .55, p = .037). Evidence further indicates that, at this early adoption stage, workflow compatibility—rather than traditional determinants—predominantly drives adoption of AI tools, suggesting that Technology Acceptance Model assumptions may require adaptation to the startup context.

**Conclusion** This review consolidates fragmented evidence into a Technology Acceptance Model-informed account of AI's role in founder decision-making, highlighting workflow compatibility as an underexamined adoption driver. Future research should examine how AI tools can complement, rather than substitute, founders' human judgment and creativity.
\end{abstract}

\section{Introduction}
One data-consistency flag before the revision: the draft states 39 studies were retained for synthesis, but the evidence split (15+7+10+11) sums to 43, not 39. I left the locked numbers untouched since they're part of your fixed SPINE, but you'll want to reconcile that arithmetic before submission.

Here is the refined introduction:

Across the United States, artificial intelligence adoption inside firms has moved from pilot project to operating infrastructure in the space of only a few years, with uptake concentrated among younger, more R&D-intensive, and digitally mature organizations (McElheran et al., 2024). Entrepreneurship scholars have been quick to register the shift: Chalmers et al. (2020) argue that AI carries implications for venture creation, resource assembly, and opportunity recognition that rival those of earlier general-purpose technologies, while Gofman and Jin (2023) show that exposure to AI curricula measurably reshapes graduates' propensity to found new ventures. Complementary streams document how AI-enabled business intelligence reconfigures managerial decision cycles (Eboigbe et al., 2023) and how the proliferation of consumer-facing generative tools is reshaping what founders and knowledge workers alike expect a decision-support "tool" to do for them (Marquis et al., 2024). Read together, this body of work suggests a field that is no longer asking whether AI matters for entrepreneurial practice, but how, where, and for whom.

Following Kaplan and Haenlein (2019), artificial intelligence is defined here as "a system's ability to interpret external data correctly, to learn from such data, and to use those learnings to achieve specific goals and tasks through flexible adaptation." This definition is broad enough to span two families of tools that founders now encounter side by side. Predictive AI—systems built to classify, forecast, or score against historical patterns—has underpinned decision-support in domains such as healthcare diagnosis and triage, where Khosravi et al. (2024) find that algorithmic recommendations can sharpen clinical judgment while still depending on human oversight for ambiguous cases. Generative AI, by contrast, produces novel text, code, or synthetic scenarios on demand; Russo (2023) documents its rapid, often improvised uptake inside software engineering teams, and Gupta (2024) provides empirical evidence that generative tools can compress the time required to move from idea to working artifact. For founders operating with limited resources, thin managerial hierarchies, and compressed timelines (Stam & Ven, 2019), the promise of both AI families appears similar: faster, better-informed choices under conditions where speed and accuracy are otherwise in tension.

Two things are therefore already well established: that AI adoption is diffusing rapidly across firms of every size (McElheran et al., 2024), and that generative and predictive tools are reshaping decision environments in adjacent domains such as software engineering, healthcare, and corporate business intelligence (Russo, 2023; Khosravi et al., 2024; Eboigbe et al., 2023). What remains largely unexamined is how these tools affect the speed and quality of decisions made by founders themselves, in the earliest, most resource-constrained phase of a venture's life, where decisions are typically made alone, under severe time pressure, and without the layered governance structures that buffer decision quality in larger organizations. The evidence that does exist is fragmented across disciplines and outcome measures—mixing case evidence, regression estimates, and survey self-reports—in ways that make it difficult to judge whether AI's effect on founder decision-making is uniformly positive, contingent, or in some settings counterproductive. Madanchian (2024) and Akinnagbe (2024) advance adjacent questions of AI adoption intention and organizational readiness, but neither synthesizes this fragmented evidence through a lens suited to founder-level decision speed and quality; to the best of our knowledge, no prior systematic review has done so.

The purpose of this review is to systematically synthesize the empirical and conceptual evidence on how AI tools affect founder decision-making speed and quality in early-stage startups, using the Technology Acceptance Model (Davis, 1989) as the organizing theoretical lens. Drawing on a corpus of 43 studies retrieved from Semantic Scholar, arXiv, and Crossref, of which 39 were retained for synthesis, the review evaluates evidence spanning regression analyses, surveys, case studies, and meta-analyses. This paper's contribution is threefold. First, it consolidates a fragmented literature—comprising 15 studies whose findings support a positive effect of AI tools on founder decision-making, 7 that report opposing or negative effects, 10 that yield mixed evidence, and 11 that bear on the question without taking a clear stance—into a single evidentiary map organized around perceived usefulness and perceived ease of use. Second, it extends the Technology Acceptance Model to the early-stage founder setting, indicating where its established predictors of adoption behave as the literature assumes and where compatibility with existing workflows, rather than usefulness or ease of use alone, may instead drive adoption decisions (Akinnagbe, 2024; Uppalapati & Nag, 2024). Third, the review moves beyond cataloguing divergent findings to articulate the boundary conditions—venture stage, founder AI literacy, task uncertainty, and tool type—that may reconcile support, opposition, and mixed results within a single contingency framework, rather than treating them as noise to be averaged away. The remainder of the paper is organized as follows: the next section details the systematic search and screening protocol; the following section presents the thematic synthesis of findings; the subsequent section develops the contingency framework; and the final section discusses theoretical and practical implications, limitations, and directions for future research.

\section{Methods}
This study employs a systematic literature review methodology following PRISMA guidelines. We searched Semantic Scholar, arXiv, and Crossref using the Boolean search string: ("AI tools") AND ("founder productivity") AND ("startup decision making") AND ("artificial intelligence"). Searches were conducted in 2026, with no lower year bound imposed, to capture the full trajectory of the field.

**PRISMA flow**: Records identified across databases: ~224; after removing duplicates: 43; screened for relevance: 43; included in synthesis: 39.

**Inclusion criteria**: (1) peer-reviewed articles or arXiv preprints with substantive empirical or theoretical content; (2) direct relevance to how do ai tools affect founder decision-making speed and quality in early-stage startups?; (3) English language. **Exclusion criteria**: abstracts with fewer than 50 words; duplicates; editorials.

Of the 39 papers included in synthesis, 38 (88%) were published within the last three years (2023–2026), confirming active research momentum. The corpus represents diverse methodological traditions including regression, survey, case study, meta-analysis.

Data extraction followed a structured, automated coding scheme applied to each paper's title and abstract, capturing: stated research questions, methodological approaches, reported findings, and explicitly signalled research gaps. Extraction and classification were performed programmatically using signal-phrase detection rather than manual coding; consequently, findings are traceable directly to the source abstracts and no inter-rater reliability statistic is reported.

Synthesis employed thematic analysis: findings were grouped into thematic clusters, frequency-weighted by citation count as a proxy for influence, and cross-validated against gap statements in each abstract. Because this review synthesizes published abstracts rather than primary data, all reported patterns should be read as characterizations of the existing literature, not as original empirical estimates.

\section{Results}
**AI as Decision-Support for Founders and Entrepreneurs (4 studies).** Across the 39 studies included in this synthesis, AI is understood broadly as a system's capacity to interpret external data, learn from it, and apply those learnings toward specific goals (Kaplan & Haenlein, 2019), and a first cluster of evidence converges on the claim that AI tools function primarily as a decision-support layer that sharpens the precision, speed, and creative range of entrepreneurial decision-making. Aziz et al. (2025) find that "AI significantly improves entrepreneurial processes by enhancing decision-making precision, facilitating opportunity recognition and fostering effective idea generation." In line with this, المعمري (2025) reports "significant effects of AI capabilities in decision making, innovation, risk mitigation and competitive advantage, all contributing to the success of entrepreneurial ventures." Indirect evidence from healthcare settings offers a partial parallel here: Scallan et al. (2026), studying clinical rather than founder populations, find that generative AI "significantly reduces the cognitive load and time required to synthesize clinical protocols" — a pattern that, if it generalizes beyond the clinical context in which it was observed, would suggest similar speed gains are available to founders, though direct founder-level replication is still needed. Consistent with this reading of AI's implications for venture creation more broadly, Chalmers et al. (2020) situate such gains within the evolving relationship between artificial intelligence and entrepreneurship. However, this decision-support narrative is not unqualified: emerging evidence indicates that while AI excels at processing information, it lacks the ability to incorporate human judgment, emotions, and creativity, which are essential in certain decision-making scenarios, suggesting boundary conditions on how far the precision and speed gains reported above can be extended.

**Adoption Drivers, Perceived Quality, and Boundary Conditions of AI Uptake (4 studies).** A second theme addresses why and how founders and professionals come to adopt AI tools in the first place, and here the evidence directly qualifies the primary theoretical lens adopted in this review, the Technology Acceptance Model. Russo (2023) finds that "at this early stage of AI integration, the compatibility of AI tools within existing development workflows predominantly drives their adoption, challenging conventional technology acceptance theories," indicating that compatibility, rather than perceived usefulness alone, may be the dominant driver in nascent adoption contexts. In line with this, Gupta (2024) identifies "social influence, domain experience, technology familiarity, system quality, training and support, interaction convenience, and anthropomorphism" as factors shaping "the pre-perception and perception phase of adoption," extending TAM's core constructs with a richer, multi-factor account of pre-adoption cognition. Consistent with a focus on perceived output quality as a precursor to continued use, Uppalapati & Nag (2024) report "a moderate positive relationship between Google Bard's accuracy and relevance (r = 0.550, p = 0.037), suggesting that as Google Bard's accuracy increases, its relevance tends to increase as well," an association consistent with the perceived-usefulness mechanism central to TAM. Marquis et al. (2024) complement this picture with a multifaceted evaluation of user perceptions of the proliferation of AI tools, reinforcing that adoption in this domain is driven by a bundle of compatibility, quality, and social factors rather than a single acceptance pathway.

**Domain Transfer: Decision-Making Evidence from Healthcare and Business Intelligence (2 studies).** Because direct, startup-specific evidence remains scarce, several included reviews draw on adjacent decision-making domains to triangulate AI's effects on decision quality. Khosravi et al. (2024) conducted "a qualitative method to conduct a systematic review of the existing reviews" with the purpose "to determine the scope of applications of AI tools in the decision-making process in healthcare service delivery networks," concluding that "AI tools have been applied in various aspects of healthcare decision-making." In line with this domain-transfer logic, Eboigbe et al. (2023) examine business intelligence transformation through AI and data analytics, pointing to a structurally similar mechanism whereby AI-enhanced data processing improves organizational decision-making outside the clinical setting. However, even in these comparatively mature domains, boundary conditions persist: Khosravi et al. (2024) note that AI's role in managing patients with multifaceted psychosocial needs remains underexplored, and that further research is needed to explore best practices and standards for implementing AI in healthcare decision-making — a caution that, by extension, applies with even greater force to the startup context, where the paucity of previous research on tool uptake from an entrepreneurial perspective highlights an urgent need for in-depth investigation.

**Structural and Human-Capital Effects of AI on Entrepreneurial Entry (2 studies).** A final theme shifts from the moment of use to the upstream question of who becomes an AI-capable founder in the first place. Gofman & Jin (2023) find that "students from the affected universities establish fewer AI startups and raise less funding," and further specify that "the brain-drain effect is significant for tenured professors, professors from top universities, and deep-learning professors," while additional evidence suggests that "unobserved city- and university-level shocks are unlikely to drive" these results, strengthening confidence that the observed effect reflects genuine human-capital reallocation rather than confounding local trends. In line with this structural view, McElheran et al. (2024) map AI adoption in America by asking who adopts, what is adopted, and where, reinforcing that access to AI expertise and tools is unevenly distributed across institutional and geographic contexts well before any founder-level decision-making effect can occur.

Read together, these four themes suggest a conceptual chain running from access to adoption to outcome. Structural human-capital dynamics (Gofman & Jin, 2023; McElheran et al., 2024) shape the initial, unevenly distributed pool of AI-capable founders; compatibility-, quality-, and social-influence-driven adoption processes (Russo, 2023; Gupta, 2024; Uppalapati & Nag, 2024; Marquis et al., 2024) determine whether and how this pool actually integrates AI into founder workflows, a process that extends and partly challenges conventional TAM assumptions; and, once adopted, AI functions as a decision-support layer that may improve the precision, speed, and creative range of entrepreneurial decision-making (Aziz et al., 2025; المعمري, 2025), with indirect evidence from healthcare and business intelligence contexts (Scallan et al., 2026; Khosravi et al., 2024; Eboigbe et al., 2023) corroborating this mechanism outside the startup setting. The gaps recurring across themes — AI's limited capacity for human judgment, emotion, and creativity in certain scenarios, its underexplored role in complex psychosocial needs, and the scarcity of dedicated research on tool uptake among founders — indicate that this chain from access to adoption to improved decision quality is neither automatic nor unconditional, underscoring the need for the startup-specific, TAM-grounded evidence that the 39 studies included in this synthesis only partially supply.

\section{Discussion}
A widespread assumption running through both practitioner commentary and parts of the academic literature is that artificial intelligence tools act as a uniform accelerant of founder cognition — that wiring generative or predictive AI into a venture's workflow mechanically produces faster, better decisions regardless of what is being decided or by whom. The 43 studies synthesized here do not support this uniformity. Instead, the corpus reveals a genuine and persistent disagreement: some studies report clear productivity and quality gains from AI-assisted decision-making, others report null effects, unintended costs, or outright caution, and a plurality occupy mixed or neutral ground. Rather than treating this split as noise to be averaged away, this discussion argues that the disagreement is itself informative: it marks the boundary conditions under which AI complements — rather than merely accelerates or destabilizes — founder judgment.

**Toward a Contingency Perspective**

Of the 43 papers in the corpus that take an identifiable stance, 15 support a positive AI–decision-making relationship, 7 oppose it, 10 report mixed results, and 11 are neutral. This is not a case where one side is simply wrong. Read closely, the supporting and opposing evidence describe different task environments, different founder populations, and different governance arrangements — and adjudicating between them, rather than reporting them side by side, is the central task of this review.

On the supporting side, gains cluster around structured, information-processing tasks. Salamzadeh et al. (2025) find that AI adoption enhances exploitative innovation, which in turn improves organizational performance — a mediated pathway consistent with AI functioning best where existing knowledge is being refined rather than novel judgment exercised. المعمري (2025) similarly frames AI's contribution as conceptual and generative, enhancing understanding of how ventures can be reorganized around AI-enabled processes. Annuš (2025) adds a human-capital qualifier within this positive cluster: founders and students with differing levels of AI knowledge use tools like ChatGPT differently, implying that the benefit is not automatic but scales with user proficiency. Gofman and Jin (2023) sharpen this further from the opposite direction — students at universities affected by a shock to AI-related human capital went on to establish fewer AI startups and raise less funding, suggesting that the positive effect is conditional on the surrounding AI-education infrastructure being intact.

On the opposing side, the concerns are not generic technology skepticism but specific to task type and governance. Davis (2024) argues for the amplified relevance of expert humans precisely because of AI, warning of the risks of omitting human judgment from decisions that carry irreducible ambiguity. Merzifonluoğlu and Gunes (2025) report that trust concerns persist and that founders and managers see a necessity in keeping final decisions under human control. George and Wooden (2023) caution that new AI-enabled paradigms carry pitfalls around quality and displacement that are easy to underestimate. Tong (2025), reviewing AI's role in product innovation, business decision-making, and marketing, likewise notes that these advantages arrive bundled with potential risks that a literature-analysis approach can surface but not resolve.

Reconciling these two bodies of evidence yields three named boundary conditions, each grounded in the studies above rather than in a simple discounting of one side.

**C1 — Task programmability.** The positive effect of AI tools on founder decision speed and quality strengthens when the decision is programmable and information-processing intensive — exploitative innovation cycles, opportunity screening, idea generation — as evidenced by Salamzadeh et al. (2025) and المعمري (2025), and weakens when the decision requires situated human judgment under ambiguity, as emphasized by Davis (2024).

**C2 — Founder AI human capital.** The positive effect strengthens with the founder's or venture team's accumulated AI-specific knowledge and tool fluency, as shown by the differentiated usage patterns in Annuš (2025), and attenuates — potentially reversing into fewer ventures founded and less capital raised — when the surrounding educational infrastructure that builds this human capital is disrupted or under-resourced, as documented by Gofman and Jin (2023).

**C3 — Governance and human-in-the-loop control.** The positive effect materializes where founders retain final decision authority and actively manage trust in the system's outputs — the condition under which Tong (2025) locates AI's advantages for product innovation, business decision-making, and marketing. Where that governance condition is absent, the effect weakens or reverses: Merzifonluoğlu and Gunes (2025) document persistent trust concerns and a founder-perceived necessity of human override, while George and Wooden (2023) show that quality and displacement pitfalls emerge precisely when such control is not exercised.

Together, C1–C3 constitute a contingency model of AI-enabled founder decision-making: AI's contribution to entrepreneurial cognition is neither uniformly positive nor uniformly overstated, but conditional on what is being decided, who is deciding, and how much control is retained over the final call.

**Theoretical Implications**

This contingency model extends the Technology Acceptance Model (Davis, 1989), the primary theoretical lens adopted in this review. TAM's original formulation explains adoption through perceived usefulness and perceived ease of use as relatively stable antecedents of behavioral intention and system use. Our findings extend this formulation by suggesting that perceived usefulness in the founder context is not a fixed property of the tool but a function of task programmability (C1): the same AI system may be perceived as highly useful for exploitative, information-intensive tasks (Salamzadeh et al., 2025) and as inadequate, even risky, for judgment-intensive decisions (Davis, 2024). This refines TAM's implicit assumption that usefulness is a general perception rather than a task-contingent one.

The founder-human-capital contingency (C2) similarly indicates that perceived ease of use — and by extension usefulness — is shaped by an antecedent TAM does not explicitly model: the depth of AI-specific human capital available to the adopter, which Gofman and Jin (2023) and Annuš (2025) suggest varies sharply across founders and educational contexts. This distinction was not explicated in TAM's original formulation, which treats users as relatively homogeneous in their capacity to form usefulness and ease-of-use beliefs.

Finally, the governance contingency (C3) challenges TAM's behavioral-intention-to-use construct as a stopping point: sustained, high-quality use appears to depend on retained human control over final decisions (Merzifonluoğlu & Gunes, 2025), an ongoing post-adoption governance dimension that sits beyond TAM's original adoption-stage focus. Collectively, these refinements suggest TAM offers a useful starting point for explaining AI adoption in founder decision-making but requires task-, capital-, and governance-contingent extensions — summarized as C1–C3 — to fully account for the mixed evidence this review has adjudicated.

**Practical Implications**

*For founders and entrepreneurs:* First, triage decisions before deploying AI — route structured, information-intensive tasks such as opportunity screening and exploitative-innovation cycles to AI tools, where the evidence of benefit is strongest (Salamzadeh et al., 2025), and route judgment-intensive or high-stakes calls (pivots, hiring, fundraising terms) through a documented human-override step (Merzifonluoğlu & Gunes, 2025; Davis, 2024). Second, budget time and money for AI literacy as an ongoing investment rather than a one-time setup cost, since differentiated knowledge levels shape how effectively tools are used (Annuš, 2025) — e.g., a recurring monthly review of which prompts, tools, and workflows are and are not working.

*For educators:* First, embed structured, credit-bearing AI training into entrepreneurship curricula rather than leaving AI fluency to incidental exposure, since disruption to university-level AI infrastructure measurably reduces subsequent AI-venture creation and funding (Gofman & Jin, 2023). Second, teach the pitfalls alongside the tools: incorporate the quality and displacement risks documented by George and Wooden (2023) directly into course modules, so literacy-building does not produce naïve overreliance.

*For policymakers:* First, treat equitable access to university-level AI education and tooling as an economic-development lever and fund it accordingly, given the venture-creation and funding gap tied to disrupted AI human-capital infrastructure (Gofman & Jin, 2023). Second, issue governance guidance for human-in-the-loop decision-making in AI-assisted ventures — for example, minimum disclosure or override standards for high-stakes, AI-informed startup decisions — in response to the trust concerns raised by Merzifonluoğlu and Gunes (2025).

*For investors:* First, do not treat AI adoption itself as a quality signal; during diligence, ask founding teams to specify which decision types they route to AI (C1) and what override mechanism governs the rest (C3). Second, look for evidence of exploitative-innovation-mediated performance gains of the kind reported by Salamzadeh et al. (2025) as a marker of effective, rather than superficial, AI integration.

**Limitations and Future Research**

This review has several limitations. First, the search was restricted to three databases — Semantic Scholar, arXiv, and Crossref — which may not capture all relevant grey literature or domain-specific outlets; however, this study can serve as a benchmark against which future reviews expanding database coverage can be calibrated. Second, the synthesis relies substantially on abstract- and sentence-level evidence rather than full-text re-analysis, which may understate methodological nuance within individual studies; this nonetheless offers a reproducible baseline for subsequent full-text coding efforts. Third, the corpus is cross-sectional, capturing a snapshot of a fast-moving literature in which 88% of papers were published in the last three years; this snapshot can itself function as a dated benchmark for tracking how the evidence balance shifts as the field matures. Fourth, publication bias may have inflated the proportion of supporting findings relative to opposing or null ones, since studies reporting clear positive effects may be more publishable; however, the presence of 7 opposing and 21 mixed-or-neutral papers in this corpus suggests the bias, if present, is only partial, and future meta-analytic work can formally quantify it. Fifth, measurement heterogeneity across the regression, survey, case-study, and meta-analysis designs identified in the corpus limits direct comparability of effect sizes; this heterogeneity nonetheless highlights where measurement harmonization would most benefit the field. Sixth, the corpus spans multiple source languages — including at least one Arabic-language study (المعمري, 2025) synthesized here via translation — introducing a risk of translation-related nuance loss that full-text, native-language re-coding in future work could resolve.

Therefore, future research should empirically test the contingency propositions articulated above (C1–C3), using designs that experimentally or longitudinally vary task programmability, founder AI literacy, and governance arrangements to isolate their moderating effects on decision speed and quality. Therefore, future research should examine human-in-the-loop governance mechanisms directly, given that best practices for maintaining human control alongside AI-assisted decision-making remain underspecified in the current corpus. Therefore, future research should investigate the pipeline between AI-focused entrepreneurship education and subsequent venture creation and funding outcomes, extending the university-level findings of Gofman and Jin (2023) into the startup context specifically, where the paucity of prior research from an entrepreneurial perspective remains a recognized gap. Therefore, future research should pursue full-text, multi-database, multilingual meta-analyses that can adjudicate the C1–C3 contingencies with greater methodological precision than the abstract-level synthesis undertaken here.

\bibliographystyle{plainnat}
\bibliography{references}

\end{document}
